\chapter*{INTRODUCCIÓN}
\thispagestyle{plain}
\addcontentsline{toc}{chapter}{INTRODUCCIÓN}

Las redes móviles de quinta generación (5G) representan un cambio paradigmático en las telecomunicaciones por su capacidad para soportar simultáneamente servicios con requisitos radicalmente heterogéneos. Esta versatilidad se materializa a través del \textit{network slicing}, capacidad habilitadora fundamental definida por 3GPP que permite crear múltiples redes lógicas independientes sobre una infraestructura física compartida, permitiendo que una misma infraestructura soporte desde \textit{streaming} de video de ultra alta definición hasta control industrial en tiempo real con latencias críticas.

La experimentación práctica con \textit{network slicing} es fundamental para la formación de ingenieros de telecomunicaciones, permitiendo comprender no solo los conceptos teóricos sino también los mecanismos técnicos que materializan la diferenciación de servicios. Sin embargo, esta experimentación enfrenta barreras significativas de accesibilidad económica y complejidad operacional. Las soluciones comerciales involucran costos prohibitivos y arquitecturas propietarias opacas, mientras que las plataformas \textit{open-source} de grado operador presentan una complejidad que excede las capacidades de laboratorios académicos típicos.

El presente trabajo propone el diseño e implementación de una plataforma de orquestación de \textit{network slicing} 5G basada en componentes \textit{open-source}, orientada específicamente a experimentación educativa y priorizando simplicidad operacional y transparencia técnica. El Capítulo 1 caracteriza esta problemática de accesibilidad y establece los objetivos, alcances y limitaciones del proyecto.